\section{The actual total cup under original local evaluation}

\begin{theorem}[The whole original total cup on the zero Cech row]
\label{thm:di-total-cup-original-zero-row}
\uses{def:di-cech-total-cup,thm:di-cech-total-cup-pure-input,thm:di-cech-total-tensor-ext}
\lean{EnvelopeProperty.DI.CechTotalCup.totalCup_zeroRow_pi_native, EnvelopeProperty.DI.CechTotalCup.totalCup_zeroRow_pi}
\leanok
Let $K$ be the original geometric relative de Rham complex, $U$ any
family of original opens, and $T$ its published Cech total complex.
For any degrees $a,b$ and original zero-row tuple, project the actual
total cup $T^a\otimes T^b\to T^{a+b}$ to the original zero Cech row
and then its original tuple factor. The entire morphism is the
original factor wedge of the two original zero-row tuple projections.
In native wedge degree the output is $b+a$; in canonical total degree
$a+b$ the formula follows by the original indexed-object equality
transport. No cover, finiteness of $U$, or comparison certificate is
required.
\end{theorem}
\begin{proof}
\leanok
The original pure-bidegree tensor inclusions jointly detect maps out
of the tensor of total terms. Test the asserted equality on bidegrees
$(q,j)$ and $(r,j')$ with $q+j=a$ and $r+j'=b$. The actual total cup
has value $(-1)^{jr}$ times the ordered cup followed by the original
inclusion of bidegree $(r+q,j+j')$. If either Cech degree is positive,
this inclusion has zero composite with the zero-row projection by
original inclusion--projection orthogonality. On the asserted right
side the corresponding input zero-row projection is zero, and
bilinearity of the original factor wedge gives zero as well.

If $j=j'=0$, the original equations give $q=a$, $r=b$ and the sign is
one. Original inclusion--projection orthogonality cancels the same
output inclusion and projection. The ordered cup projection is the
original prefix/suffix factor wedge. In Cech degree zero these two
tuples are the same original tuple and their factor restrictions are
the identity. This proves the whole native-degree equality on every
pure input, hence on the whole original tensor. Transport of the
original form index from $b+a$ to $a+b$ commutes with the original total
and product projections, giving the canonical-degree formula.
\end{proof}

\begin{theorem}[The actual original member-factor evaluation preserves the cup]
\label{thm:di-total-cup-original-factor-evaluation}
\uses{thm:di-total-cup-original-zero-row}
\lean{EnvelopeProperty.DI.CechTotalCup.totalCup_evaluation}
\leanok
For every original family member $U_i$, the native total evaluation
$T\to j_*j^*K$ takes the actual total cup to the original factor wedge
of the two native input evaluations, with the original degree
transport from $b+a$ to $a+b$.
\end{theorem}
\begin{proof}
\leanok
The native evaluation component is the original zero-row projection,
the product projection at the constant one-vertex tuple, and the
original factor restriction from that tuple intersection to $U_i$.
Apply the preceding whole-map identity to that same tuple. The
original factor restriction commutes with the original factor wedge,
as follows from restriction compatibility of the original form wedge.
Original degree transport commutes with these same factor maps.
The resulting two input composites are exactly the original native
evaluation components, proving the asserted entire morphism equality.
\end{proof}

\begin{definition}[Original member-section evaluation]
\label{def:di-total-cup-original-member-section-evaluation}
\uses{thm:di-total-cup-original-factor-evaluation}
\lean{EnvelopeProperty.DI.CechTotalCup.memberEvaluation}
\leanok
On original sections over $U_i$, compose the native factor evaluation
with the inverse of the original restriction--pushforward unit on
that same open. This gives a linear map
$\varepsilon_i^a:\Gamma(U_i,T^a)\to\Gamma(U_i,K^a)$ over the original
ring $\Gamma(U_i,\mathcal O_Y)$. The unit component is invertible:
its original cut open is $U_i\cap U_i=U_i$.
\end{definition}
\begin{proof}
\leanok
The original unit component is an original module morphism. Its
underlying map is restriction of the original sheaf to its same open,
with the actual cut-open equality. The restriction map is bijective,
so this linear map is an isomorphism. Its inverse and the native factor
evaluation therefore compose to the asserted original linear map.
\end{proof}

\begin{theorem}[Original member sections retain native local evaluation and wedge]
\label{thm:di-total-cup-original-native-local-evaluation}
\uses{def:di-total-cup-original-member-section-evaluation,thm:di-total-cup-original-factor-evaluation}
\lean{EnvelopeProperty.DI.CechTotalCup.memberEvaluation_totalCupSection, EnvelopeProperty.DI.CechTotalCup.originalLocalSection, EnvelopeProperty.DI.CechTotalCup.memberEvaluation_native}
\leanok
The map $\varepsilon_i^a$ is the original native local evaluation,
read on the original member sections through the original
open-restriction section identifications. For all whole original
sections $x\in\Gamma(U_i,T^a)$ and $y\in\Gamma(U_i,T^b)$,
\[
 \varepsilon_i^{a+b}(\mu(x,y))
 =\operatorname{cast}_{b+a=a+b}
   \bigl(\varepsilon_i^a(x)\wedge\varepsilon_i^b(y)\bigr).
\]
No exactness of sections or quasi-isomorphism hypothesis is assumed.
\end{theorem}
\begin{proof}
\leanok
The native local evaluation is the restriction of the original factor
evaluation followed by the original adjunction counit. The original
open-section identifications are natural in every original module
map. The adjunction triangle says that restricting the original unit
and following by this counit is the identity. Consequently reading
this native local evaluation on original sections is exactly the
native factor evaluation followed by the inverse original unit
component, which is $\varepsilon_i^a$.

Apply the preceding whole factor-evaluation cup identity to the
original tensor section of $x$ and $y$. The original unit commutes
with the original form wedge by restriction compatibility, and with
the original form-degree transport by naturality in that index.
The images of the two member evaluations under their original units
are the two native factor evaluations. Thus the two sides of the
displayed identity have the same image under the original invertible
target unit component. Cancel that actual unit to obtain the equality
on the whole original section modules.
\end{proof}

\begin{corollary}[Original affine readback of an evaluated global map]
\label{cor:di-total-cup-member-affine-readback}
\uses{thm:di-total-cup-original-native-local-evaluation,lem:di-affine-local-map-spec-readback}
\lean{EnvelopeProperty.DI.CechTotalCup.memberEvaluation_sectionReadback}
\leanok
Suppose $U_i$ is an original affine open and $\phi:M\to T^a$ is any
original global module map. The whole linear map obtained by applying
$\phi$ on original sections over $U_i$ and then $\varepsilon_i^a$
equals the original affine section readback of $j^*\phi$ followed by
the literal native local-evaluation component. Neither source
quasicoherence, closedness, nor a desired class comparison is assumed.
\end{corollary}
\begin{proof}
\leanok
Original affine readback is composition of original inverse-chart
restriction and genuine spectrum sections, conjugated by the original
restriction-composition and module-sections units. It therefore
preserves composition and identity. Naturality of those same original
units identifies readback of $j^*\psi$ with the original section map
of every original global morphism $\psi$.

Apply readback to the original adjunction triangle for the target
module. It identifies readback of the original counit as the inverse
of the original unit component on $U_i$; that component is invertible
by the original same-open cut equality. The literal native local
evaluation is restriction of the native factor evaluation followed
by this original counit. Composition preservation and the preceding
global-map naturality now give precisely the original section map
of $\phi$ followed by the native factor evaluation and inverse
original unit, which is $\varepsilon_i^a$.
\end{proof}
